\section{Troubleshooting}
\label{sec:Troubleshooting}

This chapter covers common failure modes, diagnostic approaches and recovery
procedures.

\subsection{Diagnostic endpoints}

Several endpoints help diagnose problems without restarting the server:

\begin{longtable}{p{54mm}p{100mm}}
\toprule
Endpoint & Purpose \\
\midrule
\code{GET /admin/overview} & Check which models exist and their health.
  Look for stale head revisions or integrity warnings. \\
\code{GET /admin/models/\{modelId\}/status} & Detailed model status:
  counts, revision, consistency. \\
\code{GET /admin/models/\{modelId\}/integrity} & Structural integrity:
  missing references, orphan nodes. \\
\code{GET /admin/auth/diagnostics} & Authorization troubleshooting:
  resolved identity, roles, mode. \\
\code{GET /admin/sessions} & Active WebSocket sessions across all
  models. Check for stale or abandoned sessions. \\
\code{GET /admin/audit/config} & Current audit configuration. \\
\bottomrule
\end{longtable}

\subsection{Common problems}

\subsubsection{Client cannot connect}

\begin{description}
\item[Symptom] Connection fails with "Failed to connect archimesh websocket".
\item[Possible causes]
  \begin{itemize}
  \item Server is not running. Check with \\
    \code{curl http://localhost:8081/admin/overview}.
  \item Wrong URL. The WebSocket URL should be \\
    \code{ws://<host>:<port>} (not \code{http}).
  \item Authorization failure. Check the auth token validity and role
    assignments. Use \code{GET /admin/auth/diagnostics}.
  \item Model does not exist. Models must be registered via
    \code{POST /admin/models/\{modelId\}} before clients can connect.
  \item Network firewall blocking WebSocket upgrade.
  \end{itemize}
\end{description}

\subsubsection{SubmitOps rejected}

\begin{description}
\item[Symptom] The server replies with a correlated \code{Error} to the client.
  The status line shows a conflict banner.
\item[Error codes]
  \begin{description}
  \item[\code{REVISION\_AHEAD}] Client's \code{baseRevision} exceeds
    server head. The client is ahead of the server --- this usually means
    the client missed a reconnect after a server state change. Solution: the
    client should resynchronize with the server. The rejection is sent to
    the submitting session and identifies its batch where available.
  \item[\code{PRECONDITION\_FAILED}] A referenced entity does not exist.
    The op references an element, relationship, view, view object or
    connection that is not in the materialized graph and was not created
    earlier in the same batch. Common causes: deleting an entity on one
    client while another client edits a diagram object representing it;
    submitting a connection create before the relationship is committed.
    The conflicting outbox entry is dropped and surfaced in the UI.
  \item[\code{LOCK\_CONFLICT}] Another client holds a notation lock on
    one of the targets. The client should retry after releasing and
    re-acquiring the lock. The conflicting outbox entry is dropped.
  \end{description}
\end{description}

\subsubsection{Materialized state corrupted}

\begin{description}
\item[Symptom] Integrity check reports missing references or orphan nodes.
\item[Cause] Bug in the apply path, manual Neo4j modification or failed
  partial writes.
\item[Solution] Run a rebuild: \\
  \code{curl -s -X POST http://localhost:8081/models/demo/rebuild}
\item[Verification] After rebuild, run integrity check again and verify
  \code{issues} is empty.
\end{description}

\subsubsection{Neo4j connection failure}

\begin{description}
\item[Symptom] Server logs show Neo4j connection errors. All operations fail.
\item[Check] Is Neo4j running? \code{docker compose ps neo4j}. Check the
  Neo4j browser at \url{http://localhost:7474}.
\item[Check] Are credentials correct? Verify \code{NEO4J\_URI},
  \code{NEO4J\_USER} and \code{NEO4J\_PASSWORD} match the running
  instance.
\item[Check] Is the Bolt port accessible? Default is 7687.
\end{description}

\subsubsection{Kafka unavailable}

\begin{description}
\item[Symptom] Server logs show Kafka publish warnings. Other clients do not
  receive real-time updates from the submitter's client.
\item[Impact] After a successful Neo4j commit, submitters can receive
  \code{OpsAccepted} while other clients await fan-out. Pending payloads
  remain durable and are retried in order after broker recovery or server
  restart. A retry may duplicate a broadcast. Pruning cannot remove pending
  publication payloads.
\item[Check] Is Kafka running? \code{docker compose ps kafka}. Check the
  Kafka UI at \url{http://localhost:8080}.
\item[Check] Do the required topics exist? Use \\
  \code{docker compose exec kafka rpk topic list --brokers localhost:9092}.
\end{description}

\subsubsection{Outbox filling up}

\begin{description}
\item[Symptom] The Archi status line shows "Outbox near capacity". Local
  edits are being queued because the client is disconnected or the
  connection is flapping.
\item[Impact] At 1000 entries, new enqueue is rejected and existing
  uncommitted work is preserved. The warning fires at 800 entries (80\%).
\item[Solution] Reconnect to the server so the outbox can drain. If
  connectivity is persistently poor, work in smaller edit batches and
  reconnect frequently.
\item[Prevention] Monitor the server's session page for clients that
  repeatedly connect and disconnect. Investigate network stability.
\end{description}

\subsubsection{Cannot overwrite a model}

\begin{description}
\item[Symptom] Import with \code{?overwrite=true} returns 409.
\item[Cause] Active sessions exist on the target model.
\item[Solution] Disconnect all clients from the model, verify with
  \code{GET /admin/sessions}, then retry. If a session is stuck, force
  delete the model and re-import.
\end{description}

\subsection{Log analysis}

The server emits structured logs. Key log categories:

\begin{longtable}{p{46mm}p{108mm}}
\toprule
Logger & Content \\
\midrule
\code{org.gautelis.archimesh} & Application-level events: submit, accept,
  reject, join, disconnect, rebuild, compact, export, import, tags. \\
\code{org.apache.kafka} & Kafka client logs (WARN by default). \\
\code{admin\_audit} (INFO) & Admin API access audit events. \\
\code{ws\_audit} (INFO) & WebSocket lifecycle audit events. \\
\bottomrule
\end{longtable}

Set \code{ARCHIMESH\_LOG\_LEVEL=DEBUG} for detailed flow tracing including
op-level detail. Note that DEBUG logging of \code{OpsBroadcast} messages
can produce very large log volumes on active models.

Log format:

\begin{lstlisting}[caption={Log format}]
yyyy-MM-dd HH:mm:ss.SSS LEVEL [logger.] (thread) message
\end{lstlisting}

Audit events include a stable JSON payload with user id, action and request
context. Parse the JSON rather than depending on free-form message text.

\subsubsection{Client-side log analysis}

The client plugin writes to the Archi workspace log
(\filepath{\textasciitilde/Archi/workspace/.metadata/.log}). Search for:

\begin{itemize}
\item \code{archimesh} --- all plugin log entries
\item \code{Connected archimesh} --- successful connections
\item \code{SubmitOps received} --- server acknowledging submits
\item \code{Queued SubmitOps} --- entries moved to outbox
\item \code{Outbox full} --- capacity reached; existing entries preserved
\item \code{Dropping stale} --- conflict resolution dropping an entry
\item \code{Error stopping Archimesh} --- shutdown errors
\end{itemize}

The \code{scripts/check-archimesh-sanity.sh} script scans Archi workspace
logs for critical sync issues, integrity warnings and connection errors.

\subsection{Recovery procedures}

\subsubsection{Repeatable process-crash drill}

Run \shellcmd{./scripts/validate-local.sh recovery} to exercise packaged server
JVMs against isolated Docker infrastructure. The drill force-stops an accepted
but unpublished batch, recovers delivery after Kafka returns, restarts after
checkpoint pruning and tail acceptance, and restores the package into an empty
database. A database uniqueness lock synchronizes a forced stop inside the
restore transaction; no partial model may remain after rollback.

The drill compares snapshots, tags, checkpoint hashes and old acceptance
ranges, then checks stale-write/delete barriers, further edits and rebuild.
It waits for the real writer lease instead of resetting it. Reports, fixture
exports and process logs are saved under
\filepath{server/target/recovery-drill/<run-uuid>/}. Ordinary gates skip this
opt-in test. Default cleanup removes only the drill's own containers, volumes
and child processes; the existing development stack is unaffected. See
\filepath{docs/recovery-drills.md} for commands, retained-infrastructure
diagnostics and the limits of the tested failure points.

\subsubsection{Pruning makes no progress}

Inspect the compact response's \code{note}, \code{checkpointRevision},
\code{watermarkRevision} and \code{eligibleBatchCount}. With no checkpoint,
all payloads remain. A recent-revision retention window can legitimately
leave no whole eligible batch. The first unpublished or unclassified legacy
batch blocks its own payload and every later one; restore Kafka delivery
before retrying. Head/commit mismatch disables pruning.

\subsubsection{Checkpoint or recovery coverage rejected}

A damaged checksum, unsupported raw state, missing receipt interval or
incomplete recovery tail is rejected. Do not bypass validation by clearing
merge metadata or changing publication flags. Restore a known-good export or
database backup. Data removed by an older server before any recovery checkpoint
cannot be reconstructed from a partial log. Rebuild validates recovery data
before replacement and rolls back transaction failures.

\subsubsection{Scheduled maintenance does not run}

It is disabled by default. Verify
\code{ARCHIMESH\_CHECKPOINT\_MAINTENANCE\_ENABLED}, the interval and revision
threshold, and the server's writer lease. The first pass waits one interval;
large catalogs rotate across passes. Search for \code{Checkpoint maintenance}
logs. Model-specific failures are isolated; later phases can be deferred even
if earlier safe phases committed. See Section~\ref{sec:ScheduledMaintenance}.

\subsubsection{Server restart}

If the server is restarted:

\begin{enumerate}
\item All WebSocket connections are closed. Clients will show
  "Disconnected" in the status line.
\item Client outboxes start queuing local edits.
\item When the server is back up, clients must manually reconnect.
\item On reconnect, the client sends a \code{Join} with its cached
  revision. The server sends a delta (if the revision gap is within the
  op-log window) or a full snapshot.
\item The outbox drains automatically after reconnect and revision hints
  are available.
\end{enumerate}

\subsubsection{Neo4j data loss}

If the Neo4j database is lost or irrecoverably corrupted:

\begin{enumerate}
\item Restore from the most recent export package using
  \code{POST /admin/models/import}.
\item If no database backup or export exists, server checkpoints are lost
  with that database. Preserve any local Archi files before reconnecting;
  they cannot reconstruct the old accepted timeline or receipts.
\item After restoring the server model, clients reconnect and receive a
  complete delta or snapshot. A new empty model does not recover the old
  timeline automatically from a client cache.
\end{enumerate}

\subsubsection{Full system recovery}

In a disaster scenario where both Neo4j and the server are lost:

\begin{enumerate}
\item Restore infrastructure: \code{docker compose up -d}.
\item Apply Neo4j schema: \code{docker compose exec -T neo4j cypher-shell ... < neo4j/schema.cypher}.
\item Start the server.
\item Import the most recent model exports.
\item Clients reconnect and receive snapshots from the restored server.
\end{enumerate}

\subsection{Performance tuning}

\begin{itemize}
\item \textbf{Neo4j heap:} The materialized graph and op-log grow with model
  activity. Monitor Neo4j heap usage and increase \code{NEO4J\_dbms\_memory\_heap\_max\_\_size}
  as needed.
\item \textbf{Kafka retention:} Long Kafka retention on ops topics consumes
  disk. If rebuild is done from Neo4j (not Kafka), set shorter retention
  for ops topics.
\item \textbf{Compaction frequency:} Run compaction periodically on active
  models after creating a verified checkpoint to reclaim published payloads.
  Review eligible and pruned batch counts; receipts and merge metadata remain.
\item \textbf{Poll interval:} The admin UI polls the server for status
  updates. Set the poll interval to 0 to disable auto-refresh if dashboard
  traffic is a concern.
\item \textbf{Consistency checks:} Keep \code{ARCHIMESH\_CONSISTENCY\_CHECK}
  disabled in production unless debugging. It adds latency to every submit.
\end{itemize}
